\section{Initialization jets and the Borel reconstruction conjecture}
\label{app:borel-population}

The population theorem above constructs the training trajectory on a fixed
positive interval.  A different question is whether a finite collection of
derivatives at initialization can approximate its predictions and loss
throughout that interval.  The Borel proposal addresses this second question.
It is an additional reconstruction conjecture; it is not a consequence of
population gradient-flow existence or uniqueness.

The precise source is Section~10 of the project's August theorem and program
note, added on 23 August 2026~\cite{ProjectBorel}.  Its original setting is
one training input and arctangent activations.  Its strongest version asks
for arbitrarily accurate finite scalar ODEs uniformly over all time, under
a positive lower bound for the kernel along the entire output interval to
the target.  The same source explicitly proposes a weaker compact-time
version.  That version is the natural companion to the positive-time theorem
in this paper.  The source's earlier uncertainty about fixed-depth flow
existence is superseded within our theorem's assumptions; its Borel
reconstruction obligations remain open.

\subsection{The observable and its initialization coefficients}

For one training input, suppress only the input index.  The physical kernel
$K(t)$ includes the fixed layer learning multipliers from the main theorem.
With physical time $t=k\eta_n$, the exact prediction equation is
\[
 \dot f(t)=2(y-f(t))K(t),\qquad \mathcal L(t)=(y-f(t))^2.
\]
There is no vanishing $\eta_n$ in this population equation.  The source uses
an older time convention with a fixed learning multiplier.  In the present
convention its normalized time is simply $\tau=2t$.

Write $f_0=f(0)$ and $e_0=y-f_0$.  Assume $e_0>0$ and $K(0)>0$; reversing
the output direction handles $e_0<0$.  If $e_0=0$, the one-input squared-loss
flow is stationary.  By continuity, $K$ remains positive on some fixed
interval $[0,T_1]$.  Direct integration of the residual equation gives
\[
 y-f(t)=e_0\exp\!\left(-2\int_0^t K(u)\,du\right)>0.
\]
Thus $f$ increases strictly there, and the initialized trajectory defines
an output-dependent kernel by
\[
 \kappa(f(t))=K(t).
\]
The scalar function $\kappa$ is different from the fixed layer learning
multipliers $\kappa_\ell$.
This is a function along the one initialized trajectory.  It does not say
that arbitrary population states with the same prediction have the same
kernel.  In particular, this scalar description does not replace the full
population state law for arbitrary restarts.

Suppose the required actual derivatives at initialization exist.  Define
\[
 d_p=\left.\frac{d^p f}{d\tau^p}\right|_{\tau=0}
     =2^{-p}f^{(p)}(0),
 \qquad
 \kappa(f_0+s)\sim\sum_{p=0}^{\infty}c_p s^p.
\]
Here $s$ is only the scalar output increment, not a hidden preactivation.
Dividing $df/d\tau$ by $y-f$ and formally inverting $f(\tau)-f_0$ determines
$c_p$ from the finite initialization jet $e_0,d_1,\ldots,d_{p+1}$.  The first
coefficients are
\begin{align*}
 c_0&=\frac{d_1}{e_0},\\
 c_1&=\frac{e_0d_2+d_1^2}{e_0^2d_1},\\
 c_2&=\frac{d_3}{2e_0d_1^2}
       -\frac{d_2^2}{2e_0d_1^3}
       +\frac{d_2}{e_0^2d_1}
       +\frac{d_1}{e_0^3}.
\end{align*}
For example, $d_1=e_0c_0$ and
$d_2=(-c_0+e_0c_1)d_1$, which gives the first two formulas.  Differentiating
once more gives the third.  The denominators are nonzero because
$d_1=e_0K(0)>0$.  A zero limiting readout does not obstruct this condition:
the readout contribution to $K(0)$ can already be positive while all hidden
initial velocities vanish.

The source proposes computing each fixed-order $d_p$ from initialization,
the architecture, and the update rule by mean-field peeling.  This involves
two obligations: computing the finite-order coefficient and showing that it
is the derivative of the actual population flow.  Uniform convergence of
network paths alone does not permit differentiating the width limit.

For orientation, a degree-$q$ raw truncation would produce the scalar ODE
\[
 \dot f_{[q]}=2(y-f_{[q]})
       \sum_{p=0}^{q}c_p(f_{[q]}-f_0)^p,
 \qquad f_{[q]}(0)=f_0.
\]
It formally matches the first $q+1$ prediction derivatives.  This is a
useful finite example, but matching more derivatives does not establish
convergence on a positive interval.  The purpose of Borel reconstruction is
to handle cases where these raw truncations fail.

\subsection{What Borel reconstruction adds to a formal jet}
\label{app:borel-conditions}

For a candidate exponent $\sigma>0$, define the transformed series
\[
 \mathcal B_\sigma\kappa(\xi)
 =\sum_{p=0}^{\infty}
   \frac{c_p}{\Gamma(1+\sigma p)}\xi^p.
\]
A coefficient bound
\[
 |c_p|\le C A^p\Gamma(1+\sigma p)
\]
makes this series converge near $\xi=0$.  Such a bound is called a Gevrey
coefficient bound.  The constants here are numerical bounds, not additional
population fields.  Growth of derivatives and growth of power-series
coefficients differ by a factor $p!$: a derivative bound of the form
$C A^p(p!)^r$ corresponds, after dividing by $p!$, to coefficient growth
of order $(p!)^{r-1}$.  In the convention of the summability literature
cited below, the summability order associated with our denominator is
$1/\sigma$.  No exponent $\sigma$ for the present neural kernel is supplied
by the population existence theorem.

The intended identity is
\begin{equation}
 \kappa(f_0+s)=\int_0^\infty e^{-u}
               \mathcal B_\sigma\kappa(su^\sigma)\,du.
 \label{eq:borel-actual-kernel}
\end{equation}
This is an identity with the actual kernel, not merely a definition of a
function having its formal coefficients.  For every $s>0$, the integral
samples arbitrarily far along the positive $\xi$-ray.  Local convergence of
the transformed series therefore leaves a continuation problem.  One also
needs growth control to make the integral exist and be uniformly controlled.
For instance, for constants $C>0$ and $b\ge0$, the sufficient bound
\[
 |\mathcal B_\sigma\kappa(\xi)|
 \le C\exp(b\xi^{1/\sigma}),\qquad \xi\ge0,
\]
gives a common integrable envelope for $0\le s\le S$ whenever
$bS^{1/\sigma}<1$.  Indeed the integrand is then bounded by
$C\exp[-(1-bS^{1/\sigma})u]$.

Finally, finitely many initialization coefficients must produce a declared
approximation to this continued transform.  Write $B_N$ for such a rational
approximation, and let $(u_j,w_j)_{j=1}^M$ be a finite quadrature rule for the
Laplace integral.  The resulting kernel and prediction equation are
\begin{align*}
 \kappa_{N,M}(f_0+s)&=\sum_{j=1}^M w_jB_N(su_j^\sigma),\\
 \dot f_{N,M}&=2(y-f_{N,M})\kappa_{N,M}(f_{N,M}),
 \qquad f_{N,M}(0)=f_0.
\end{align*}
All model-dependent coefficients must come from initialization and the
declared reconstruction rule.  Positive-time samples or fitted unknown
trajectory data cannot be inserted into $B_N$.  The desired sufficient
hypothesis is uniform convergence to the actual kernel on the output
interval being used:
\[
 \sup_{v\in I}|\kappa_{N,M}(v)-\kappa(v)|\longrightarrow0.
\]
Here $I$ is a fixed interval of reached outputs.  The finite orders may grow
with the requested accuracy, but are chosen independently of width and of
the vanishing GD step.  Pad\'e continuation is one possible choice; its
convergence and control of poles must be proved.  Appendix~\ref{app:pade-borel}
gives sufficient conditions and alternatives.

The mechanism for transferring transform convergence through the integral
is elementary.  Cut the integral at a large finite $U$.  Up to $U$, only
the compact transform interval $[0,SU^\sigma]$ is needed.  Uniform
approximation there controls this part.  A common integrable envelope
controls the remaining tail, and finite quadrature controls the compact
integral.  A cutoff construction can also be used when only the true
transform has a tail bound; this is explained in Appendix~\ref{app:pade-borel}.

It is essential that this is a reconstruction procedure beyond raw
truncation.  Exact integration of a truncated Borel power series gives
\[
 \int_0^\infty e^{-u}
   \sum_{p=0}^N\frac{c_p(su^\sigma)^p}{\Gamma(1+\sigma p)}\,du
 =\sum_{p=0}^N c_p s^p.
\]
The equality follows from the defining Gamma integral.  Thus dividing by
Gamma factors and then integrating the polynomial simply returns the
possibly divergent Taylor truncation.

\subsection{Why existence and uniqueness do not identify the sum}

The July monograph records the surviving obstruction: smooth functions can
have the same complete initialization jet and differ at every positive
time~\cite[Section~5.7.6]{ProjectMonograph}.  Consider the smooth function
\[
 g(x)=
 \begin{cases}
  \exp(-1/x^2),&x>0,\\
  0,&x\le0.
 \end{cases}
\]
Every derivative of $g$ at zero vanishes, and $g'$ is bounded.  Hence the ODE
\[
 \dot x=1+g(x),\qquad x(0)=0
\]
has a unique global solution.  Differentiating its equation at zero gives
$x'(0)=1$ and $x^{(p)}(0)=0$ for $p\ge2$.  Its formal Taylor series is
therefore just $t$, which is convergent and Borel summable.  Yet
\[
 x(t)=t+\int_0^t g(x(u))\,du>t\qquad(t>0).
\]
The candidate sum $t$ fails the ODE by $g(t)$.  This is a logical example,
not a counterexample within the neural model.  It explains why one cannot
invoke uniqueness merely because a proposed reconstruction has matching
initial derivatives: uniqueness applies to solutions of the equation.

There are two direct ways to complete identification.  One can prove
\eqref{eq:borel-actual-kernel} for the actual observable, using suitable
analytic or quasianalytic assumptions.  Alternatively, one can reconstruct
the full population state, verify that its sum solves the same integral
GF equations in the strong solution class, and then apply uniqueness.
Matching only the prediction or kernel jet does not verify those full
state equations.  Sectorial quasianalytic classes provide relevant analytic
criteria, but their hypotheses are additional to real-time GF
existence~\cite{LastraMalekSanz}.

The theorem in the main text permits $C^{1,1}$ activations.  It does not
assume all-order differentiability, and its moment and stability estimates
do not establish all-order jets, Gevrey growth, or quasianalyticity.  Even
$C^\infty$ smoothness alone would not remove the flat-function obstruction.
Arctangent remains a natural smoother test case for this separate program.

The old reconstruction audit considered proving the width limit through
Borel summation, which would require common finite-width continuation and
identification estimates.  That route is unnecessary here.  The population
flow and width/GD limit have already been obtained independently.  It is
enough to identify a reconstruction of the population observable and then
use the triangle inequality at the end of this appendix.

\subsection{Uniform kernel error gives uniform scalar training error}
\label{app:borel-scalar-error}

Here is the precise local implication, including the time margin.  Choose
$0<T<T_1\le T_*$ such that $K>0$ on $[0,T_1]$, and set
$I=[f_0,f(T_1)]$.  Suppose a continuous approximate kernel $\kappa_N$ satisfies
\[
 \sup_I|\kappa_N-\kappa|\le\varepsilon<\min_I\kappa.
\]
The single index $N$ now denotes any chosen finite reconstruction, including
its quadrature.  Define $f_N$ by
\[
 \dot f_N=2(y-f_N)\kappa_N(f_N),\qquad f_N(0)=f_0.
\]
At each output in $I$, its velocity is between
$1-\varepsilon/\min_I\kappa$ and $1+\varepsilon/\min_I\kappa$ times the
true velocity.  This comparison can be proved without a derivative bound
for $\kappa$: the time to reach output $v$ is exactly
\[
 \int_{f_0}^{v}\frac{dw}{2(y-w)\kappa(w)},
\]
and the approximate clock has $\kappa_N$ in its denominator.  Both clocks
are strictly increasing, so inversion gives
\[
 f\!\left(\left(1-\frac{\varepsilon}{\min_I\kappa}\right)t\right)
 \le f_N(t)\le
 f\!\left(\left(1+\frac{\varepsilon}{\min_I\kappa}\right)t\right).
\]
The comparison is valid on $[0,T]$ provided
$(1+\varepsilon/\min_I\kappa)T\le T_1$.  This margin avoids evaluating the
exact trajectory beyond its proved interval.  Positive velocity also gives
existence and uniqueness by the same clock inversion on the range used.
Since $|\dot f|\le2e_0\max_I\kappa$ and
$|\dot{\mathcal L}|\le4e_0^2\max_I\kappa$, we obtain
\begin{align*}
 \sup_{t\le T}|f_N(t)-f(t)|
 &\le 2e_0T\frac{\max_I\kappa}{\min_I\kappa}\,\varepsilon,\\
 \sup_{t\le T}|\mathcal L_N(t)-\mathcal L(t)|
 &\le 4e_0^2T\frac{\max_I\kappa}{\min_I\kappa}\,\varepsilon.
\end{align*}
The positive lower bound here follows locally from $K(0)>0$ and continuity.
It is not a lower bound all the way to the target.

For completeness, the source's stronger all-time implication also has a
short proof.  Suppose instead that $\kappa$ and $\kappa_N$ are continuous
on the full output interval $I=[f_0,y]$, with the same positive minimum and
error condition.  The bounded positive kernels give global scalar flows
approaching $y$.  The clock comparison now holds for all time.  Moreover,
\[
 \mathcal L(t)\le e_0^2e^{-4t\min_I\kappa},
 \qquad
 \sup_{t>0}t|\dot{\mathcal L}(t)|
 \le\frac{e_0^2\max_I\kappa}{\mathrm e\min_I\kappa}.
\]
The second inequality maximizes
$4t e_0^2\max_I\kappa\,e^{-4t\min_I\kappa}$.  Integrating the derivative
with respect to $\log t$ between the compared clocks yields
\[
 \sup_{t\ge0}|\mathcal L_N(t)-\mathcal L(t)|
 \le
 \frac{e_0^2\max_I\kappa}{\mathrm e\min_I\kappa}
 \left[-\log\!\left(1-\frac{\varepsilon}{\min_I\kappa}\right)\right].
\]
Here $\mathrm e$ is Euler's number.  This verifies the source's conditional
all-time bound.  Our local GF theorem does not supply its full-interval
hypotheses or prove that the neural kernel has such a global extension.

\subsection{The additional reconstruction needed for several inputs}
\label{app:borel-matrix-error}

For the fixed dataset in the main theorem,
\[
 \dot f_a(t)=-2\sum_b\omega_bK_{ab}(t)(f_b(t)-y_b).
\]
An individual prediction can turn, or remain still while other predictions
and hidden fields evolve.  The one-input output coordinate therefore does
not automatically extend to this system.

A precise additional conjecture would reconstruct the matrix $K(t)$ from
its physical-time initialization derivatives.  Suppose this produces a
continuous, initialization-only matrix $K_N(t)$ on $[0,T]$.  Let
$\Omega=\operatorname{diag}(\omega_1,\ldots,\omega_m)$ and assume
\[
 \varepsilon_N=
 \sup_{t\le T}
 \left\|\Omega^{1/2}(K_N(t)-K(t))\Omega^{1/2}\right\|_{\operatorname{op}}
 \longrightarrow0.
\]
Define the finite prediction approximation by
\[
 \dot f_{N,a}(t)=-2\sum_b\omega_b(K_N(t))_{ab}(f_{N,b}(t)-y_b),
 \qquad f_{N,a}(0)=f_a(0),
\]
and put $\mathcal L_N(t)=\sum_a\omega_a(f_{N,a}(t)-y_a)^2$.
Vectors $f_N,f,y$ below collect these $m$ scalar entries.  This finite
known-coefficient equation is different from the source's output-only
scalar ODE.  It can be made autonomous by adding one clock coordinate; this
introduces no unknown trajectory data if $K_N$ is actually computed from
initialization.

The true weighted kernel $\Omega^{1/2}K(t)\Omega^{1/2}$ is positive
semidefinite.  The symmetric part of its approximation is therefore bounded
below by $-\varepsilon_N$ times the identity.  Differentiating the weighted
residual energy gives
\[
 \mathcal L(t)\le\mathcal L(0),\qquad
 \mathcal L_N(t)\le e^{4\varepsilon_Nt}\mathcal L(0).
\]
To estimate the prediction difference, subtract the two vector equations:
\begin{align*}
 \frac{d}{dt}\bigl[\Omega^{1/2}(f_N-f)\bigr]
 ={}&-2\Omega^{1/2}K\Omega^{1/2}
                  \bigl[\Omega^{1/2}(f_N-f)\bigr]\\
    &-2\Omega^{1/2}(K_N-K)\Omega^{1/2}
                  \bigl[\Omega^{1/2}(f_N-y)\bigr].
\end{align*}
The first term cannot increase the Euclidean norm.  The second has norm at
most $2\varepsilon_Ne^{2\varepsilon_Nt}\sqrt{\mathcal L(0)}$.
Integration from the zero initial difference proves
\begin{align*}
 \sup_{t\le T}\|\Omega^{1/2}(f_N(t)-f(t))\|_2
 &\le (e^{2\varepsilon_NT}-1)\sqrt{\mathcal L(0)},\\
 \sup_{t\le T}|\mathcal L_N(t)-\mathcal L(t)|
 &\le (e^{4\varepsilon_NT}-1)\mathcal L(0).
\end{align*}
For the loss bound use the difference of squared norms and the two residual
norm bounds.  This argument does not require kernels at different times to
commute or a positive minimum kernel eigenvalue.  If $K_N$ is also positive
semidefinite, both residual norms are at most $\sqrt{\mathcal L(0)}$, and
the respective bounds improve to
$2\varepsilon_NT\sqrt{\mathcal L(0)}$ and
$4\varepsilon_NT\mathcal L(0)$.

Uniform entrywise error at most $\varepsilon$ suffices for
$\varepsilon_N\le\varepsilon$: the squared Frobenius norm of the weighted
matrix difference is at most
$\varepsilon^2\sum_{a,b}\omega_a\omega_b=\varepsilon^2$.
Thus an entrywise Borel reconstruction, with identification and uniform
error control, would be enough.  Neither the source's one-input conjecture
nor the present GF theorem establishes that extra matrix reconstruction.

\subsection{Consequences for networks and the exact time quantifier}

Fix a desired accuracy and a finite initialization-based approximation with
population prediction error at most $\varepsilon$ on $[0,T]$.  In the
multiple-input case, the main theorem and the triangle inequality give
\begin{align*}
 &\sup_{t\le T}\|\Omega^{1/2}(f_n^{\mathrm{GD}}(t)-f_N(t))\|_2\\
 &\qquad\le
 \sup_{t\le T}\|\Omega^{1/2}(f_n^{\mathrm{GD}}(t)-f(t))\|_2
 +\varepsilon
 =o_{\mathbb P}(1)+\varepsilon.
\end{align*}
For one input, replace the weighted norm by absolute value.  Loss follows
from the same bounded-residual estimates.  A finite amount of initialization
information would therefore suffice for a prescribed prediction and loss
accuracy throughout a fixed positive interval, independently of width and
the vanishing step size.  The approximation order is chosen for accuracy
before taking the width limit.  Arbitrary joint growth of that order with
width would require an additional quantitative estimate.  Existence of a
sufficient order also does not give a certified stopping rule unless a
computable error bound is available.

To cover every $t<T_*$, identified reconstruction must hold on every compact
$[0,T]$ with $T<T_*$, or on all of $[0,T_*]$.  Reconstruction on one smaller
fixed interval $[0,T_B]$, with $T_B>0$, gives a conclusion only there.  If
that interval shrinks to zero with width, step size, or approximation order,
it establishes nothing at any fixed positive time.  GF existence alone
does not extend the reconstruction range of an initialization expansion.

Finally, this conjecture concerns specified prediction and kernel
observables.  It does not automatically reconstruct hidden field laws,
hidden Gram matrices, or population operators; each needs its own identified
reconstruction or a theorem for the full state.  It does not remove nonlinear
feature learning, give a finite-dimensional law for every possible
population restart, or establish an all-time result for the broad
$C^{1,1}$ class.  The Stieltjes conditions in the next appendix are useful
sufficient analytic assumptions, not assumptions of the original Borel
conjecture and not consequences of the GF proof.
